\section{Related work}
\label{sec:related-work}

Our comparison separates the security of a ratcheting construction, the
exposure represented by its state, and the interpretation of that state in a
durable implementation.  This places the reservation problem alongside
existing cryptographic and systems results without identifying a new storage
write, a completed exchange, and a recovered secret as the same event.

\paragraph{Ratcheting and exposure of future messages.}
Cohn-Gordon et al. analyze Signal's X3DH and Double Ratchet core as a multi-stage
authenticated key-exchange protocol \cite{signal2020}.  Their model represents
the branching structure of ratchet stages and specifies admissible key
disclosures for a challenged stage.  This establishes an important precedent
for deriving a precise cryptographic object from a deployed state machine.
The object studied here includes a contribution's recoverable randomness
before its KEM computation, together with the record versions that retain it.
Our witnesses connect these earlier objects to later roots and application
messages.

Auerbach et al. introduce SCKA and vulnerable epoch and message sets
\cite{scka2025}.  Definition~3.2 includes current chains, already produced
future-epoch keys, and pending out-of-order message keys.  Theorem~3.3 gives a
secure-messaging construction from secure SCKA, an appropriate KDF, and AEAD.
Their Section~1.2 explicitly identifies state loss and session handling as a
separate setting.  We use their set-based perspective and investigate how
original reservations and persistent-record acquisition determine exposure at
that representation boundary.  In particular, our delayed-receipt witness is
compared with the state of both directions and their retained contributions,
rather than assigned a universal bound from a confirmed-epoch counter.
% Primary comparison: USENIX Security 2025 paper, Sections 1.2 and 3.1--3.3,
% especially Definition 3.2 and Theorem 3.3.  Future-epoch exposure and
% direction/index-sensitive vulnerable sets are already present there.

The timing of secret generation is itself an established design consideration.
Section~8.5 of Signal's Double Ratchet specification discusses deferring ratchet
key generation until it is needed for sending, thereby shortening a key's
lifetime \cite{doubleratchet2025}.  Our contribution concerns the operational
meaning of that lifetime when a durable reservation fixes the key before the
generation call runs.  An implementation can defer computation while retaining
the original coins.  The relevant freshness event is then the reservation,
and its identity must survive retries and subsequent image writes.

\paragraph{Session handling and state-loss resilience.}
Cremers, Jacomme, and Naska distinguish individual ratcheting sessions from the
conversation assembled by a messaging application \cite{sessionhandling2023}.
Their Signal experiments and Tamarin model show how session selection and
retry behavior permit a clone to remain active, and they analyze two mechanisms
for stronger conversation guarantees.  This is directly relevant to the scope
of our disclosure interfaces: a session label identifies a protocol instance,
whereas shared authentication or wrapping keys can connect several instances.
The Continuity model follows actual key identities, and the aggregate-copy
execution examines a concrete storage relationship between two sessions.

Cremers, Medinger, and Naska subsequently prove incompatibilities between full
conversation PCS and specified state-loss resilience requirements
\cite{stateloss2025}.  Their Definitions~1 and~2 require continued communication
after dynamic or static state loss; the distinction includes losing evolving
communication state and losing persistent identity state.  Our crash/resume
setting retains a bound journal and resumes its original operations under
explicit authentication assumptions.  The recovery theorem states exactly
which reservation, identity, and record-acquisition histories are admissible.
Thus, the retained state and independent authentication in our claim identify
the particular recovery setting being studied.  The delayed-receipt and
aggregate-copy witnesses explain how that setting behaves during honest
asynchronous progress.
% Primary comparison: Cremers--Medinger--Naska, Sections 3.1--3.3,
% Definitions 1--2, and the stated impossibility results.
% Snapshot disclosure, recovery from a retained journal, and loss of all
% dynamic/static state must retain their separate meanings.

\paragraph{Robustness on unreliable networks.}
Fischlin, G\"unther, and Janson parameterize channel correctness by supported
delivery histories and develop corresponding robustness notions
\cite{robustchannels2024}.  Their analysis of QUIC and DTLS~1.3 accounts for
loss, reordering, replay policies, and packet-number reconstruction.  In
particular, adversarial traffic should preserve acceptance of supported honest
ciphertexts.  This provides a useful methodological parallel: legitimate
network behavior belongs inside a security model's operational semantics.
We retain delayed receipts and asymmetric cutover when reasoning about
reservation provenance.  Delivery admissibility and secret freshness are
separate predicates on the same execution; synchronizing both peers before
allowing the next reservation would change the execution that our witnesses
and theorem analyze.
% Primary comparison: Robust Channels, Definitions 3.1--3.2, Section 4,
% and the QUIC/DTLS analyses in Sections 6--7.

\paragraph{Post-quantum ratchets and communication comparisons.}
Dodis et al.'s Triple Ratchet combines classical and post-quantum continuous
key agreement, using erasure coding to distribute the larger post-quantum
messages \cite{tripleratchet2025}.  Its Section~4 develops the hybrid
construction and security argument.  Remark~4.3 already accounts for exposure
of next-epoch material sampled while its transmission is in progress.
We therefore compare original-reservation persistence with an existing
analysis of retained future material.  Our SPQR and whole-KEM experiments
report the scheduling and retention consequences of particular implementations;
their authentication mechanisms, control opportunities, and persistence
interfaces remain explicit experimental parameters.
% Primary comparison: extended Triple Ratchet paper, Sections 4.1--4.2
% and Remark 4.3; Section 5 defines its RKEM abstraction.

The reference corpus separates the pinned upstream SPQR code from the
experimental whole-KEM construction.  The latter is a comparison control with
fixed local-send intervals.  It contributes an executable full-message
exchange for the same schedule corpus.  A case reports the concrete message
keys derived from one endpoint snapshot and the public transcript; its
\texttt{not\_derived} label records the output of that derivation program.
This convention preserves the distinction between an observed exposure set
and the security statement supplied by a cryptographic theorem.  It also
allows reordered traffic to retain its direction and message index instead of
being compressed into a single recovery time.

\paragraph{Mechanized ratchet analyses.}
Cheval, Jacomme, and Richards give an automated analysis of the Double Ratchet
that includes out-of-order messages, key-compromise scenarios, and PCS
\cite{doubleratchet2026}.  Their work also identifies and analyzes fixes for
forward-secrecy defects in the specification and implementation.  These results
make detailed ratchet state and compromise interfaces central comparison
points for our symbolic analysis.

The CCS~2026 program includes Bhargavan et al.'s analysis of SPQR and its
implementation \cite{spqr2026}.  The authors' published account describes
ProVerif models used during protocol design and hax/F* verification of the
Rust core's absence of panics and selected correctness properties
\cite{spqrverification2025}.  Our comparison with this work concerns the state
represented by an exposure and the provenance preserved by durable replay.
The present contribution combines a checked recovery condition with concrete
reservation-disclosure executions and an explicit normalization of record
aliases.
% Metadata for the forthcoming CCS 2026 paper is verified against the
% official program.  The implementation-verification details above are
% attributed specifically to the author account, not inferred from its title.

H\'erouard et al. mechanize a computational PCS proof for an asymmetric
ratchet in Squirrel \cite{squirrel2026}.  Their analysis covers two parties
and unbounded ratchet steps with a fixed, non-adaptive compromise scenario.
The proof distinguishes non-deduction of intermediate roots from
real-or-random security of derived keys, and develops rules for reasoning about
these properties in the random-oracle model.  This distinction informs our
computational interface: each intermediate-root claim is attached to the
view in which it is used, and message-key indistinguishability is a separate
conclusion.  A convention that samples ideal names in advance is also distinct
from giving an adversary a persistent token from which actual operation coins
can be recovered; the disclosure interface determines which values enter its
view.
% Primary comparison: Squirrel paper, Section 4.2 (health/compromise
% predicates), Section 4.3 (security proof), and Appendix E (root secrecy).

\paragraph{Key derivation with an exposed predecessor.}
Continuity's root extraction invokes HMAC with the previous root in the key
position and the new contribution in the message position.  RFC~5869 specifies
this argument order for HKDF-Extract \cite{rfc5869}.  When the predecessor root
is exposed, the relevant secret therefore occurs in HMAC's message input.
Backendal et al. characterize the input-domain feasibility conditions for
HMAC's swapped PRF use and give security bounds under compression-function
assumptions \cite{hmacswap2023}.  These results supply a precise reference for
that extraction step.  Our computational accounting additionally identifies
the record interface, source attribution, and joint root/traffic experiment
to which an instantiation must apply.

Across these comparisons, the specific object added by the present study is
the provenance of retained operation randomness through reservation, record
acquisition, replay, and asynchronous confirmation.  The normalization result
preserves that provenance at the record interface, the symbolic theorem gives
a sufficient recovery condition, and the native witnesses show why the
condition refers to the original events.
